\section{Sliding Window OR}
\label{sec:cses-sw-or}

\problemstatement{Given an array of $n$ integers and a window size $k$,
output the bitwise OR of every contiguous window of size $k$.}

\begin{patternbox}
OR is \textbf{not invertible}: removing an element cannot simply undo its
bits, because another element in the window may still set them. The fix is
to track, for each bit position, \emph{how many} window elements set that
bit -- then the OR has a bit set exactly when its count is positive.
\end{patternbox}

\subsection*{Per-bit counts}

Maintain an array $\textit{cnt}[b]$ = number of elements in the current
window whose bit $b$ is set. Adding an element increments $\textit{cnt}[b]$
for each of its set bits; removing an element decrements them. The OR of
the window has bit $b$ set iff $\textit{cnt}[b]>0$. Reconstruct the OR by
summing $2^b$ over all bits with positive count. This restores $O(1)$
amortised updates (constant number of bits) despite OR's lack of an
inverse.

\begin{ideabox}[Count Contributors per Bit]
When an aggregate is not invertible, replace ``does any element set this
bit?'' with ``how many do?'' A count is invertible (increment/decrement),
and the OR is recovered by testing each count against zero. The same idea
handles AND (count zeros) and other non-invertible bitwise folds.
\end{ideabox}

\subsection*{Algorithm}

\begin{enumerate}
  \item For the first window, increment $\textit{cnt}[b]$ for every set
        bit of each element.
  \item To slide: add the new element's bits, remove the old element's
        bits.
  \item Output $\sum_{b:\,\textit{cnt}[b]>0}2^b$ for each window.
\end{enumerate}

\subsection*{Dry run}

\dryrun{$a=[1,2,3]$ (bits: $1{=}001$, $2{=}010$, $3{=}011$), $k=2$.}

\begin{itemize}
  \item Window $[1,2]$: bit$0$ count $1$, bit$1$ count $1$ $\Rightarrow$
        OR $=3$.
  \item Slide to $[2,3]$: remove $1$ (bit$0\to0$), add $3$ (bit$0\to1$,
        bit$1\to2$) $\Rightarrow$ bits $0,1$ set, OR $=3$.
  \item Output $3,3$.
\end{itemize}

\subsection*{C++ implementation}

\begin{lstlisting}
#include <bits/stdc++.h>
using namespace std;

int main() {
    int n, k;
    cin >> n >> k;
    vector<int> a(n);
    for (int& x : a) cin >> x;

    const int BITS = 31;
    vector<int> cnt(BITS, 0);
    auto add = [&](int x, int delta) {
        for (int b = 0; b < BITS; b++) if (x & (1 << b)) cnt[b] += delta;
    };
    auto currentOR = [&]() {
        int res = 0;
        for (int b = 0; b < BITS; b++) if (cnt[b] > 0) res |= (1 << b);
        return res;
    };

    for (int i = 0; i < k; i++) add(a[i], +1);
    cout << currentOR() << " ";
    for (int i = k; i < n; i++) {
        add(a[i], +1); add(a[i - k], -1);         // add new, remove old
        cout << currentOR() << " ";
    }
    cout << "\n";
    return 0;
}
\end{lstlisting}

\begin{complexitybox}
\textbf{Time Complexity.} $O(n\cdot\text{BITS})$ -- constant $31$ bits per
element. \textbf{Space Complexity.} $O(\text{BITS})$.
\end{complexitybox}

\begin{takeawaybox}
Non-invertible aggregates like OR are made sliding-friendly by counting
contributors per bit, so removal becomes a decrement. This
``count-then-threshold'' conversion is the general recipe whenever an
operation lacks an inverse.
\end{takeawaybox}
